\chapter{Covariant fields and Gaussian elimination}
\label{cbc:chap:covariant}

A contraction of complexes need not commute with moving the
insertion points. The abelian line algebra gives a simple example:
its cohomology consists of charges, but translating a positive mode
can produce a charge. Differential forms on the space of positions
record the missing homotopy. We construct that homotopy and retain
the logarithmic connection that survives the contraction.

In Gaussian elimination, the hidden quadratic form can vary with
the fields that remain. Its normalization then contributes to their
effective differential. The lifts of their antifields also change.
We compute both effects in a finite polynomial Batalin--Vilkovisky
complex. Products of lifts require curvature corrections, whose
connected contractions determine higher comparison maps.

Throughout, $k$ is a field of characteristic zero and differentials
have cohomological degree one. Odd differentiation is on the left.
Tensor products and operator extensions use the Koszul rule.
The line model and the finite Gaussian model have the separate
coefficient rings and completions specified below.

\section{Charges and spectral translation}
\label{cbc:cf:line}

Fix $\kappa\in k^\times$. Consider the graded commutative algebra
\[
 Y_\kappa=k[B_0,B_1,\ldots]\otimes
                 \Lambda(C_0,C_1,\ldots),
 \qquad |B_n|=0,\quad |C_n|=1.
\]
Define derivations $Q,T,L$ of degrees $1,0,-1$ by
\begin{align}
 QB_n&=\kappa nC_{n-1},& QC_n&=0,\notag\\
 TB_n&=nB_{n-1},& TC_n&=nC_{n-1},\notag\\
 LC_n&=\kappa^{-1}B_n,& LB_n&=0.
 \label{cbc:cf:operators}
\end{align}
Terms with index $-1$ are zero. The relations
\[
 Q^2=L^2=0,\qquad QL+LQ=T,\qquad [T,L]=[T,Q]=0
\]
follow on generators and hence on the algebra.

Give $k[J]$ degree zero and zero differential. Let
$\pi:Y_\kappa\to k[J]$ send $B_0$ to $J$ and every
other generator to zero. Let $\iota(J)=B_0$.

Let $N$ count the $C_n$ and the $B_n$ with $n>0$,
with multiplicity. Define
\begin{equation}\label{cbc:cf:line-h}
 h=N^{-1}\sum_{n\geq0}
       \frac{B_{n+1}}{\kappa(n+1)}\partial_{C_n}
 \quad\text{on }N>0,\qquad h=0\quad\text{on }N=0.
\end{equation}

The sum is finite on each polynomial. If its numerator is $H$,
then $QH+HQ=N$: each pair $(B_{n+1},C_n)$ contributes
its Euler operator. Also $H^2=0$, and $H$ preserves $N$.
It follows that
\begin{equation}\label{cbc:cf:line-sdr}
 Qh+hQ=1-\iota\pi,\qquad
 h^2=h\iota=\pi h=0,\qquad \pi\iota=1.
\end{equation}
Thus $(\pi,\iota,h)$ is a strong deformation retract,
linear over $k[B_0]$. Its projection fails to preserve
translation already on $B_1$, since $\pi(TB_1)=J$.

\begin{proposition}[The de Rham correction]\label{cbc:cf:cartan}
Let $\epsilon$ have degree one, with $\epsilon^2=d\epsilon=0$.
Place forms on the left and put $W=1+\epsilon L$.
Then
\[
 (Q+\epsilon T)W=WQ.
\]
The dg algebra map
\[
 \widetilde\pi=\pi(1-\epsilon L):
 \bigl(\Lambda(\epsilon)\otimes Y_\kappa,Q+\epsilon T\bigr)
 \longrightarrow
 \bigl(\Lambda(\epsilon)\otimes k[J],0\bigr)
\]
is a quasi-isomorphism. Together with $\iota$ and
$WhW^{-1}$ it is a strong deformation retract.
\end{proposition}
\begin{proof}
The operator $\epsilon L$ is a degree-zero derivation with square
zero. Its exponential is $W$, whose inverse is $1-\epsilon L$.
Extending $Q$ over forms gives $Q(\epsilon a)=-\epsilon Qa$.
Consequently the coefficient of $\epsilon$ in the intertwining
equation is $T-QL=LQ$. Conjugating
\eqref{cbc:cf:line-sdr} proves the retract identities.
Since $L\iota=0$, the conjugated inclusion is still $\iota$.
\end{proof}

For example,
$\widetilde\pi(C_0)=-\epsilon J/\kappa$.
The one-form term records the translation homotopy of a state
whose strict charge projection vanishes.

\begin{proposition}[Polynomial spectral transport]\label{cbc:cf:spectral}
On polynomial inputs, the maps
\[
 E_z=(1+dz\,L)e^{zT}:
 Y_\kappa\longrightarrow
 \Omega^\bullet(\mathbb A^1_z)\otimes Y_\kappa
\]
are dg algebra maps and satisfy $E_{z+w}=E_zE_w$ after
pullback to two parameters. Let $\Delta_0$ make every $B_n,C_n$
primitive, and set
$\widetilde\Delta_z=(E_z\otimes1)\Delta_0$.
The spectral coproducts satisfy
\[
 (\widetilde\Delta_u\otimes1)\widetilde\Delta_v
       =(1\otimes\widetilde\Delta_v)\widetilde\Delta_{u+v}.
\]
On a generator $a$ their common value is
\[
 E_{u+v}a\otimes1\otimes1+
 1\otimes E_va\otimes1+1\otimes1\otimes a.
\]
\end{proposition}
\begin{proof}
The derivation $T$ lowers the mode index, so $e^{zT}$ is finite
on each polynomial. The identity $QL+LQ=T$, with forms on the
left, gives $(Q+d_z)E_z=E_zQ$. Both factors of $E_z$ are
algebra automorphisms over forms. They commute under composition
because $[T,L]=0$. The possible product of two one-form terms
vanishes by $L^2=0$. This proves the group law.

Each generator is sent to a linear combination of generators.
The primitive coproduct therefore gives the displayed equality,
which extends to all inputs by multiplicativity.
\end{proof}

The charge projection of this family is
\begin{align*}
 (\pi\otimes\pi)\widetilde\Delta_z(B_0)&=J_1+J_2,\\
 (\pi\otimes\pi)\widetilde\Delta_z(B_n)&=z^nJ_1
                                      &&(n>0),\\
 (\pi\otimes\pi)\widetilde\Delta_z(C_n)&=
                  \kappa^{-1}z^nJ_1\,dz &&(n\geq0).
\end{align*}
Pullback along $z\mapsto tz$, including
$dz\mapsto t\,dz+z\,dt$, gives a polynomial homotopy
to the primitive charge map at $t=0$. Its composition law
is the group law just proved.

\section{The completed line connection}
\label{cbc:cf:completion}

To sum all modes, fix a finite ordered set $I$ of positions and
write
\[
 R_I=k[z_i\ (i\in I),z_{ij}^{-1}\ (i<j)],
 \qquad z_{ij}=z_i-z_j,\qquad
 Y_I=\bigotimes_{i\in I}Y_\kappa.
\]

Give $B_n$ weight $n$ and $C_n$ weight $n+1$. We use the
product completion
\begin{equation}\label{cbc:cf:jet-ring}
 \widehat Y_{I,\mathrm{jet}}
       =\prod_{w\geq0}(R_I\otimes Y_I)_w.
\end{equation}
At each weight the coefficients are polynomials in the finitely
many zero modes $B_0^{(i)}$. Localization in position differences
precedes completion, so the finite pole order may vary with $w$.
Multiplication involves only finitely many summands at any fixed
weight. Differential forms have weight zero.

For two positions put $z=z_1-z_2$ and define
\begin{align}
 r(z)&=\sum_{n,m\geq0}(-1)^n\binom{n+m}{n}
       \frac{C_n\otimes B_m-B_n\otimes C_m}{z^{n+m+1}},
       \label{cbc:cf:r}\\
 \mathcal A&=r(z)+dz\,L_1r(z),&
 \rho&=\kappa^{-1}J_1J_2\,\frac{dz}{z},
       \label{cbc:cf:connection}
\end{align}
where $L_1=L\otimes1$. Let $D_0=Q+d_{\mathrm{dR}}$.
The sums define continuous multiplication operators on the
form-valued version of \eqref{cbc:cf:jet-ring}.

The positive-weight element
\begin{equation}\label{cbc:cf:primitive}
 b(z)=-\frac1\kappa\sum_{N\geq1}\frac{z^{-N}}N
             \sum_{a+c=N}(-1)^a\binom Na B_a\otimes B_c
\end{equation}
is a primitive for all but the charge residue.

\begin{lemma}\label{cbc:cf:primitive-identity}
In the completed coefficient algebra,
\[
 Qb=r,\qquad
 d_{\mathrm{dR}}b=dz\,L_1r-\iota\rho,\qquad
 D_0b=\mathcal A-\iota\rho.
\]
In particular $D_0\mathcal A=0$ and
$(D_0+\mathcal A\cdot)^2=0$ on coefficient modules.
\end{lemma}
\begin{proof}
Applying $Q$ to $B_a$ in \eqref{cbc:cf:primitive} gives,
for $a\geq1$, the coefficient
\[
 -\frac{(-1)^aa}{N}\binom Na
       =(-1)^{a-1}\binom{N-1}{a-1}
\]
of $C_{a-1}\otimes B_c$. Applying $Q$ to $B_c$ gives
$-(-1)^a\binom{N-1}{a}$ on $B_a\otimes C_{c-1}$.
These are exactly the terms of \eqref{cbc:cf:r}.

Differentiating $-z^{-N}/N$ gives $z^{-N-1}dz$.
The resulting sum is $dz\,L_1r$ with its $N=0$ term
$\kappa^{-1}B_0\otimes B_0\,dz/z$ removed.
The first two identities follow weight by weight.
The form $\rho$ is closed, so the last identity implies
$D_0\mathcal A=0$. Since $\mathcal A$ has degree one
in a graded commutative algebra, $\mathcal A^2=0$.
\end{proof}

For instance the weight-one term is
\[
 b(z)=\frac{B_1\otimes B_0-B_0\otimes B_1}{\kappa z}
                  +\text{terms of weight at least two}.
\]
It contains the first-mode information that appeared in the
spectral projection.

On several sites use the sum of the numerators in
\eqref{cbc:cf:line-h}, divided by the total number of ghosts and
positive-index $B$'s. This homotopy and $Q$ preserve weight,
so their contraction extends continuously to
\eqref{cbc:cf:jet-ring}.

\begin{theorem}[The line module contraction]\label{cbc:cf:line-retract}
Let $G=e^{-b}$, which is defined weight by weight. The maps
\[
 I_{\mathcal A}=G\iota,\qquad
 P_{\mathcal A}=\pi,\qquad
 H_{\mathcal A}=GhG^{-1}
\]
give a continuous strong deformation retract
\[
 \left(\prod_{w\geq0}
       \Omega^\bullet(R_{\{1,2\}})\otimes(Y_{\{1,2\}})_w,
       D_0+\mathcal A\cdot\right)
 \ \rightleftarrows\
 \left(\Omega^\bullet(R_{\{1,2\}})\otimes k[J_1,J_2],
       d_{\mathrm{dR}}+\rho\cdot\right).
\]
These are maps of module complexes.
\end{theorem}
\begin{proof}
Positive weight makes both $G$ and $G^{-1}$ well defined.
Lemma~\ref{cbc:cf:primitive-identity} gives
\[
 (D_0+\mathcal A)G=G(D_0+\iota\rho).
\]

The original homotopy is linear over the zero-mode polynomial
ring. Its extension over forms anticommutes with both
$d_{\mathrm{dR}}$ and multiplication by the odd form $\rho$.
Thus \eqref{cbc:cf:line-sdr} also contracts the $\rho$-twisted
complexes. Conjugation gives the asserted homotopy and its side
conditions. Every term of $b$ contains a positive-index $B$,
so $\pi G=\pi G^{-1}=\pi$ and $P_{\mathcal A}I_{\mathcal A}=1$.
All identities are finite at each fixed weight.
\end{proof}

The endpoint vector is $I_{\mathcal A}(1)=G$.
On the charge complex the differential sends $1$ to $\rho$.
A further scalar gauge with constant term one in the charges
would require, at the coefficient of $J_1J_2$, a rational
primitive of $-\kappa^{-1}dz/z$. Every rational derivative
has zero residue, so such a gauge cannot remove this connection
over $R_{\{1,2\}}$. Analytic choices of a logarithm would change
the coefficient domain.

\begin{corollary}\label{cbc:cf:all-positions}
For any finite $I$, set
\[
 b_I=\sum_{i<j}b_{ij},\qquad
 \mathcal A_I=\sum_{i<j}(r_{ij}+dz_{ij}L_ir_{ij}),\qquad
 \rho_I=\kappa^{-1}\sum_{i<j}J_iJ_j\,d\log z_{ij}.
\]
Theorem~\ref{cbc:cf:line-retract} holds with
$G_I=e^{-b_I}$ and the total homotopy.
Conjugating specified continuous coefficient maps by these
gauges preserves their composition identities.
\end{corollary}
\begin{proof}
The primitive identities add over pairs and all exponent terms
commute. For a map $f$ from a source model to a target model,
the conjugate is $G_{\mathrm t}^{-1}fG_{\mathrm s}$.
The adjacent gauge factors cancel in every composite.
Splitting a set of positions into groups splits $b_I$ into
internal and cross terms, whose exponentials multiply.
\end{proof}

At three positions the pair sum includes $13$ with its spectator.
At four positions it includes the crossed pairs. The corollary
uses the distinct-position ring $R_I$ and presupposes the
coefficient maps being conjugated. Extension through a collision
requires further maps.

The completion also restricts spectral transport. The sequence
$B_n$ tends to zero in jet weight, but for fixed $z\in k^\times$
\[
 \pi(e^{zT}B_n)=z^nJ
\]
does not tend to zero in the discrete weight-zero coefficient
ring. Numerical translation therefore has no continuous extension
to \eqref{cbc:cf:jet-ring}. Proposition~\ref{cbc:cf:spectral}
has polynomial inputs. Giving a formal translation parameter
weight one defines a different, combined completion.

\section{A Gaussian complex with retained coordinates}
\label{cbc:cf:gaussian}

The line contraction leaves a connection on its charge module.
We now compute the corresponding effect of eliminating finitely
many hidden fields. The hidden quadratic form may depend on the
retained coordinates, so its normalization contributes to their
effective differential.

Let $R$ be a commutative $k$-algebra with commuting derivations
$\partial_a$, $1\leq a\leq r$. We write them as derivatives
in retained coordinates $y_a$. Let $A$ be an invertible symmetric
$d\times d$ matrix over $R$, put $C=A^{-1}$, and choose $U\in R$.
The ring contains $(\det A)^{-1}$ and is preserved by the
derivations.

Introduce hidden even coordinates $x_1,\ldots,x_d$,
a degree-zero parameter $\hbar$, and odd antifields
$\eta_{x_i},\eta_{y_a}$ of degree $-1$. The algebras are
\[
 \mathcal B=R[x_1,\ldots,x_d,\hbar]\otimes\Lambda(\eta_x,\eta_y),
 \qquad
 \mathcal B_{\mathrm{ret}}=R[\hbar]\otimes\Lambda(\eta_y).
\]
All hidden-coordinate dependence is polynomial.

For $S=U+x^{\mathsf t}Ax/2$, the quantum differential is
\begin{equation}\label{cbc:cf:Q}
 Q_S=\sum_i\bigl((Ax)_i-\hbar\partial_{x_i}\bigr)
                                      \partial_{\eta_{x_i}}
       +\sum_a\bigl(\partial_aS-\hbar\partial_a\bigr)
                                      \partial_{\eta_{y_a}}.
\end{equation}
Let $q$ range over the hidden and retained coordinates.
The even operators $\partial_qS-\hbar\partial_q$ commute,
and the odd derivatives anticommute. Hence $Q_S^2=0$.
This is a complex over $k[\hbar]$, since the retained
derivatives need not be $R$-linear.

For example, take $R=k[y,(1+y^2)^{-1}]$ and
$S=y^2/2+(1+y^2)x^2/2$. Then
\[
 Q_S\eta_y=y+yx^2.
\]
Eliminating $x$ must replace the hidden quadratic term and lift
$\eta_y$ to an element whose differential has no hidden variable.

The differential has order two for the polynomial product.
For homogeneous $f,g$, denote its product deviation by
\begin{equation}\label{cbc:cf:deviation}
 B_S(f,g)=Q_S(fg)-(Q_Sf)g-(-1)^{|f|}f(Q_Sg).
\end{equation}
It is divisible by $\hbar$. This failure of the Leibniz
identity will determine the corrections to products of lifts.

\subsection{The finite heat contraction}

Let $N_x$ count hidden polynomial degree and hidden antifield
number. With repeated hidden indices summed, put
\[
 d_{\mathrm{cl}}=(Ax)_i\partial_{\eta_{x_i}},\qquad
 h_{\mathrm{cl}}=N_x^{-1}\eta_x^{\mathsf t}C\partial_x
                   \quad(N_x>0),
\]
and set $h_{\mathrm{cl}}=0$ on hidden constants.
Let $\epsilon_x$ set $x=\eta_x=0$, and let $i$ include
$\mathcal B_{\mathrm{ret}}$. The Euler identity gives
$d_{\mathrm{cl}}h_{\mathrm{cl}}+
h_{\mathrm{cl}}d_{\mathrm{cl}}=1-i\epsilon_x$.
The other contraction identities follow as in
\eqref{cbc:cf:line-sdr}.

Define
\begin{equation}\label{cbc:cf:heat}
 W_x=\exp\left(-\frac{\hbar}{2}
                    \partial_x^{\mathsf t}C\partial_x\right),
 \qquad P=\epsilon_xW_x^{-1},\qquad
 h_x=W_xh_{\mathrm{cl}}W_x^{-1}.
\end{equation}
The exponentials terminate on each polynomial. Their commutator
with multiplication by $Ax$ gives
$Q_x=W_xd_{\mathrm{cl}}W_x^{-1}$, where $Q_x$ is the hidden
part of \eqref{cbc:cf:Q}. Thus $(P,i,h_x)$ is a strong
deformation retract for $Q_x$. The projection $P$ is Wick
evaluation with covariance $\hbar C$ and kills every hidden
antifield.

\begin{theorem}[Variable Gaussian contraction]\label{cbc:cf:gaussian-sdr}
Let $\delta_y=Q_S-Q_x$. The finite formulas
\begin{equation}\label{cbc:cf:gaussian-maps}
 I_A=(1+h_x\delta_y)^{-1}i,\qquad
 H_A=h_x(1+\delta_yh_x)^{-1},\qquad P_A=P
\end{equation}
give a strong deformation retract of $(\mathcal B,Q_S)$ onto
$\mathcal B_{\mathrm{ret}}$ with differential
\begin{equation}\label{cbc:cf:effective}
 Q_A=\sum_a(E_a-\hbar\partial_a)\partial_{\eta_{y_a}},
 \qquad
 E_a=\partial_aU+\frac{\hbar}{2}\operatorname{tr}(C\partial_aA).
\end{equation}
Every map has polynomial $\hbar$-coefficients.
\end{theorem}
\begin{proof}
The operator $\delta_y$ lowers retained antifield number by one,
whereas $h_x$ preserves it. The inverse series in
\eqref{cbc:cf:gaussian-maps} terminate after at most $r$
perturbations on each input.

To verify the contraction, set
$t=(1+\delta_yh_x)^{-1}\delta_y
  =\delta_y(1+h_x\delta_y)^{-1}$.
The identity
\[
 Q_xt+tQ_x=-t\,iP\,t
\]
follows from $(Q_x+\delta_y)^2=0$ and
$Q_xh_x+h_xQ_x=1-iP$.
Indeed multiply on the left by $1+\delta_yh_x$ and on
the right by $1+h_x\delta_y$. The difference becomes
\[
 Q_x\delta_y+\delta_yQ_x+
 \delta_y(h_xQ_x+Q_xh_x+iP)\delta_y=0.
\]
Substitution in
\[
 I_A=i-h_xti,\qquad P_A=P-Pth_x,\qquad
 H_A=h_x-h_xth_x,\qquad Q_A=Pti
\]
gives the chain equations and
\[
 Q_SH_A+H_AQ_S=1-I_AP_A,\qquad
 P_AI_A=1,\qquad H_A^2=H_AI_A=P_AH_A=0.
\]

Every value of $h_x$ has a hidden antifield, which $\delta_y$
cannot remove. Hence $P\delta_yh_x=0$, so $P_A=P$.
This also kills every higher term of the transferred
differential. Its only hidden quadratic expectation is
\[
 P\left(\frac12x^{\mathsf t}(\partial_aA)x\right)
             =\frac{\hbar}{2}\operatorname{tr}(C\partial_aA).
\]
This proves \eqref{cbc:cf:effective} on every retained
antifield degree. The finite heat and perturbation formulas
give the stated coefficient ring.
\end{proof}

In the one-variable example, the formulas give
\[
 E_y=y+\frac{\hbar y}{1+y^2},\qquad
 I_A(\eta_y)=\eta_y-\frac{xy}{1+y^2}\eta_x.
\]
The extra antifield term cancels the hidden part of the
differential and leaves the Gaussian variance in the force.

The effective force is the closed algebraic one-form
\[
 dU+\frac{\hbar}{2}\operatorname{tr}(A^{-1}dA).
\]
Here $d\log\det A$ abbreviates
$\operatorname{tr}(A^{-1}dA)$. No logarithm need belong to $R$.
The Gaussian normalization can be retained as the rank-one
module with connection
$d+\tfrac12\operatorname{tr}(A^{-1}dA)$.
For $A=y$ on $\mathbb G_m$, a local analytic horizontal
section is $y^{-1/2}$, with sign monodromy around zero.
A positive real quadratic form with a positive branch supplies
an additional scalar state.

\section{Curvature and products of antifields}
\label{cbc:cf:curvature}

The finite inclusion first sends a retained antifield to
\[
 I_A(\eta_{y_a})=\eta_{y_a}-\eta_x^{\mathsf t}M_ax,
 \qquad M_a=\frac12C\partial_aA.
\]
The relation $\partial_aA=M_a^{\mathsf t}A+AM_a$ cancels
the hidden quadratic force. Products of these unary lifts
have a further defect, measured by the curvature of $M$.
We allow any connection satisfying the same metric relation.

Choose matrices $M_a$ over $R$ with
\begin{equation}\label{cbc:cf:metric}
 \partial_aA=M_a^{\mathsf t}A+AM_a,\qquad
 F_{ab}=\partial_aM_b-\partial_bM_a+[M_a,M_b].
\end{equation}
The canonical choice $C\partial_aA/2$ always exists.
Metric compatibility implies
$F_{ab}^{\mathsf t}A+AF_{ab}=0$, so $F_{ab}C$
is antisymmetric.

Define
\begin{equation}\label{cbc:cf:hats}
 \widehat\eta_a=\eta_{y_a}-\eta_x^{\mathsf t}M_ax,\qquad
 R_{ab}=\frac{\hbar}{2}\eta_x^{\mathsf t}F_{ab}C\eta_x.
\end{equation}
The element $R_{ab}=-R_{ba}$ has degree $-2$. Taking the
trace of \eqref{cbc:cf:metric} gives
$E_a=\partial_aU+\hbar\operatorname{tr}M_a$.
Thus the effective differential is independent of this choice
of metric connection.

\begin{lemma}[The pair correction]\label{cbc:cf:pair}
The elements in \eqref{cbc:cf:hats} satisfy
\begin{align}
 Q_S\widehat\eta_a&=E_a,&
 B_S(\widehat\eta_a,\widehat\eta_b)
       &=\hbar\eta_x^{\mathsf t}F_{ab}x,\notag\\
 Q_SR_{ab}&=-\hbar\eta_x^{\mathsf t}F_{ab}x,&
 B_S(R_{ab},R_{cd})&=0.
 \label{cbc:cf:pair-identities}
\end{align}
Hence $\widehat\eta_a\widehat\eta_b+R_{ab}$ is a chain
lift of $\eta_{y_a}\eta_{y_b}$ compatible with the unary
lifts. For the canonical connection,
$F_{ab}=-[C\partial_aA,C\partial_bA]/4$.
\end{lemma}
\begin{proof}
In $Q_S\widehat\eta_a$, metric compatibility cancels the
quadratic terms. The quantum derivative of
$-\eta_x^{\mathsf t}M_ax$ contributes
$\hbar\operatorname{tr}M_a$. The horizontal vector fields
associated to $\widehat\eta_a$ are
$\partial_a-(M_ax)\cdot\partial_x$. Their commutator is
$-(F_{ab}x)\cdot\partial_x$. Expansion of
\eqref{cbc:cf:deviation} gives the second identity.

For an antisymmetric matrix $K$, the classical hidden
differential sends $\eta_x^{\mathsf t}K\eta_x/2$ to
$-\eta_x^{\mathsf t}KAx$. Apply this with $K=F_{ab}C$.
Its quantum hidden derivative vanishes because $K$ has no
$x$-dependence, and its retained differential vanishes because
it has no retained antifield. This proves the third identity.
Two curvature pairings have neither hidden-coordinate dependence
nor retained antifields, so their product deviation is zero.

Finally $\theta=C\,dA$ satisfies
$d\theta+\theta\wedge\theta=0$.
The curvature of $\theta/2$ is $-\theta\wedge\theta/4$,
which proves the canonical formula.
\end{proof}

For $f\in R$, the additional term of the differential on a
product by $f$ is
$-\hbar\sum_a(\partial_af)\partial_{\eta_{y_a}}$.
The pair correction has no retained antifield, so the lift
also obeys the chain equation on $f\eta_{y_a}\eta_{y_b}$.

\begin{example}\label{cbc:cf:curved-example}
For retained coordinates $y,z$, take
\[
 A=\begin{pmatrix}1+y^2&yz\\yz&1+z^2\end{pmatrix},
 \qquad U=\frac12(y^2+z^2),\qquad
 R=k[y,z,(1+y^2+z^2)^{-1}].
\]
The canonical correction and effective coefficients are
\[
 R_{yz}=-\frac{\hbar(y^2+z^2)}
                   {2(1+y^2+z^2)^2}\eta_{x_1}\eta_{x_2},
 \qquad
 E_y=y+\frac{\hbar y}{1+y^2+z^2},\quad
 E_z=z+\frac{\hbar z}{1+y^2+z^2}.
\]
At $(y,z)=(1,0)$ the curvature is
\[
 F_{yz}=\begin{pmatrix}0&-1/8\\1/4&0\end{pmatrix},
 \qquad R_{yz}=-\frac{\hbar}{8}\eta_{x_1}\eta_{x_2}.
\]
Thus the product of the two specified unary lifts needs a
nonzero correction.
\end{example}

\section{The inclusion in every antifield degree}
\label{cbc:cf:matching-section}

Four retained antifields can require two pair corrections at once.
A sum over partial matchings gives the full inclusion and explains
the role of the Bianchi identity.

For an increasing list $L=(a_1,\ldots,a_m)$, set
$J_M(\varnothing)=1$ and define
\begin{equation}\label{cbc:cf:matching}
 J_M(L)=\widehat\eta_{a_1}J_M(a_2,\ldots,a_m)
 +\sum_{j=2}^m(-1)^jR_{a_1a_j}
             J_M(a_2,\ldots,\widehat{a_j},\ldots,a_m).
\end{equation}
Let $I_M(\eta_{y_{a_1}}\cdots\eta_{y_{a_m}})=J_M(L)$
and extend by $R[\hbar]$-linearity. Each unmatched label
contributes $\widehat\eta_a$, each pair contributes $R_{ab}$,
and collecting the labels gives the exterior sign. In degree
four the term with two pairs is
\[
 R_{ab}R_{cd}-R_{ac}R_{bd}+R_{ad}R_{bc},
 \qquad a<b<c<d.
\]

\begin{theorem}[The matching inclusion]\label{cbc:cf:matching-theorem}
The map
$I_M:(\mathcal B_{\mathrm{ret}},Q_A)\to(\mathcal B,Q_S)$
is unital and commutes with the differentials. The Wick
projection $P$ of \eqref{cbc:cf:heat} satisfies $PI_M=1$.
All coefficients lie in $R[\hbar]$.
\end{theorem}
\begin{proof}
Put $\nabla_aF_{bc}=\partial_aF_{bc}+[M_a,F_{bc}]$.
Besides \eqref{cbc:cf:pair-identities}, direct differentiation
gives
\begin{equation}\label{cbc:cf:bianchi}
 B_S(\widehat\eta_a,R_{bc})
   =-\frac{\hbar^2}{2}
       \eta_x^{\mathsf t}(\nabla_aF_{bc})C\eta_x.
\end{equation}
Indeed, differentiate the pairing
$\eta_x^{\mathsf t}F_{bc}C\eta_x$ and use
$\partial_aC=-M_aC-CM_a^{\mathsf t}$.
The derivative of $F$ and the two actions on its normal
factors combine into $\nabla_aF$. The cyclic sum of
\eqref{cbc:cf:bianchi} is zero by the Bianchi identity
$\nabla_aF_{bc}+\nabla_bF_{ca}+\nabla_cF_{ab}=0$.

Apply the second-order product formula to each matching.
Differentiating an unmatched factor gives its $E_a$ term.
The deviation of two unmatched factors cancels the differential
of the matching in which they are paired, by
\eqref{cbc:cf:pair-identities}. Deviations between an unmatched
factor and a pair cancel in triples by \eqref{cbc:cf:bianchi}.
Pair--pair deviations vanish.

These are all the terms because
$Q_S$ has order two. The sign $(-1)^j$ in
\eqref{cbc:cf:matching} moves label $a_j$ past the $j-2$
intermediate odd labels to form the pair. Thus each cancellation
has the exterior sign of the differential on the retained
monomial.

For coefficients $f\in R$, the additional derivative is
$-\hbar\sum_a(\partial_af)\partial_{\eta_{y_a}}$.
In the matching formula, retained odd differentiation removes
the unmatched factor with that label. It therefore commutes with $I_M$.
This proves the complete chain equation over $k[\hbar]$.

The map $P$ is a chain map by
Theorem~\ref{cbc:cf:gaussian-sdr}. Every term in $I_M$ other
than the original retained monomial contains a hidden antifield,
so $P$ kills it. Hence $PI_M=1$. Each matched edge contributes
one factor of $\hbar$, and there are finitely many matchings.
\end{proof}

\begin{example}[Several simultaneous pairs]\label{cbc:cf:rotations}
Take $A=I$ and hidden and retained ranks $2\ell$. Let $J_a$
have entries $1,-1$ at $(2a-1,2a),(2a,2a-1)$ and zero
elsewhere. Choose
\[
 M_{2a-1}=0,\qquad M_{2a}=y_{2a-1}J_a,
 \qquad 1\leq a\leq\ell.
\]
Distinct rotations commute. The only nonzero pairings are
$R_{2a-1,2a}=\hbar\eta_{x_{2a-1}}\eta_{x_{2a}}$.
The all-paired term in
$I_M(\eta_{y_1}\cdots\eta_{y_{2\ell}})$ is
\[
 \hbar^\ell\eta_{x_1}\cdots\eta_{x_{2\ell}}\ne0.
\]

Thus these chosen connections require arbitrarily many pairs
as the rank increases. At fixed hidden rank, the exterior
algebra bounds the number of pairs. The same constant metric
also admits $M=0$, whose inclusion is the ordinary one.
\end{example}

\subsection{A homotopy and the flatness criterion}

The matching formula need not coincide with the particular
inclusion $I_A$ chosen by finite perturbation. The same
projection permits a direct comparison of their homotopies.

\begin{proposition}\label{cbc:cf:matching-homotopy}
Set $H_M=H_A(1-I_MP)$. Then
\[
 Q_SH_M+H_MQ_S=1-I_MP,\qquad
 H_M^2=H_MI_M=PH_M=0.
\]
For two metric connections $M,M'$, the degree-minus-one map
$K_{M,M'}=H_A(I_M-I_{M'})$ satisfies
\[
 Q_SK_{M,M'}+K_{M,M'}Q_A=I_M-I_{M'}.
\]
\end{proposition}
\begin{proof}
Since $1-I_MP$ is a chain map, its commutator with $Q_S$
after multiplication by $H_A$ is
\[
 (1-I_AP)(1-I_MP)=1-I_MP,
\]
using $PI_M=1$. The equations $PH_M=H_MI_M=0$ follow
from the same identity and $PH_A=0$. In the expansion of
$H_M^2$, every term contains $H_A^2$ or $PH_A$ and vanishes.
Finally apply $Q_SH_A+H_AQ_S=1-I_AP$ to $I_M-I_{M'}$.
Its projection is zero, giving the comparison homotopy.
\end{proof}

\begin{proposition}\label{cbc:cf:flat}
The matching inclusion $I_M$ is a homomorphism of graded
commutative algebras if and only if $F_{ab}=0$ for all $a,b$.
A flat metric connection exists over $R$ whenever there is
$G\in\operatorname{GL}_d(R)$ for which $G^{\mathsf t}AG$
is diagonal and invertible.
\end{proposition}
\begin{proof}
Multiplicativity on two retained antifields requires $R_{ab}=0$.
The matrix $F_{ab}C$ is antisymmetric. Independence of the
exterior monomials and the absence of $\hbar$-torsion make this
equivalent to $F_{ab}C=0$, hence to $F_{ab}=0$.
Conversely, zero curvature removes every matched term, leaving
the algebra map $\eta_{y_a}\mapsto\widehat\eta_a$.

Write $D=G^{\mathsf t}AG=\operatorname{diag}(d_1,\ldots,d_d)$.
In this frame the matrices
\[
 N_a=\frac12\operatorname{diag}
       (d_1^{-1}\partial_ad_1,\ldots,d_d^{-1}\partial_ad_d)
\]
are metric compatible. Their commutators vanish and their
logarithmic one-forms are closed, so they are flat. Differentiating
$D=G^{\mathsf t}AG$ shows that
$M_a=GN_aG^{-1}-(\partial_aG)G^{-1}$ is compatible with $A$.
Its curvature is the conjugate of that of $N$, hence zero.
\end{proof}

For the metric of Example~\ref{cbc:cf:curved-example},
\[
 G=\begin{pmatrix}1&-yz/(1+y^2)\\0&1\end{pmatrix}
 \quad\text{gives}\quad
 G^{\mathsf t}AG=
 \operatorname{diag}\left(1+y^2,\frac{1+y^2+z^2}{1+y^2}\right).
\]
This gives a multiplicative inclusion over
$k[y,z,(1+y^2)^{-1},(1+y^2+z^2)^{-1}]$.
It uses an additional inverse beyond the original ring.

\begin{problem}[A flat metric connection]\label{cbc:cf:flat-question}
Does the metric
$\left(\begin{smallmatrix}1+y^2&yz\\yz&1+z^2\end{smallmatrix}\right)$
admit a flat metric connection over
$k[y,z,(1+y^2+z^2)^{-1}]$?
\end{problem}

The nonzero canonical curvature only obstructs multiplicativity
of its specified lifts. Changing a connection can change its
curvature. Changing its frame conjugates curvature and cannot
make a nonzero curvature vanish.

\section{Frames and Gaussian normalization}
\label{cbc:cf:frames}

A change of hidden coordinates also changes the retained
antifields. The density determines the corresponding quantum
correction. For $x=G(y)x'$ with
$G\in\operatorname{GL}_d(R)$, use
\begin{align}
 A'&=G^{\mathsf t}AG,&
 M'_a&=G^{-1}M_aG+G^{-1}\partial_aG,\notag\\
 \eta_x&=G^{-\mathsf t}\eta_x',&
 \eta_{y_a}&=\eta_{y_a}'-
       \eta_x'{}^{\mathsf t}G^{-1}(\partial_aG)x'.
 \label{cbc:cf:frame-change}
\end{align}
These are the cotangent transformation rules. For a density
factor $\nu\in R^\times$, define
\[
 Q_\nu=Q_S-\hbar\sum_a\nu^{-1}(\partial_a\nu)
                                     \partial_{\eta_{y_a}}.
\]
In a retained coordinate chart this is the differential for
the volume density $\nu(y)\,dx_1\cdots dx_d\,dy_1\cdots dy_r$.

\begin{proposition}\label{cbc:cf:density}
The matching inclusion is preserved by
\eqref{cbc:cf:frame-change}. For $Q_\nu$, the retained force is
\begin{equation}\label{cbc:cf:density-force}
 dU+\frac{\hbar}{2}\operatorname{tr}(A^{-1}dA)
                                    -\hbar\nu^{-1}d\nu.
\end{equation}
It is invariant under $\nu'=(\det G)\nu$.
\end{proposition}
\begin{proof}
Substitution gives $\widehat\eta_a=\widehat\eta'_a$.
Also $F'_{ab}=G^{-1}F_{ab}G$, so
$\eta_x^{\mathsf t}F_{ab}C\eta_x$ is unchanged.
Every matching is therefore preserved.

Both $I_M$ and $P$ commute with retained odd differentiation.
The density operator subtracts the last term of
\eqref{cbc:cf:density-force} on the retained complex.
Finally $\det A'=(\det G)^2\det A$.
Its logarithmic derivative change cancels that of the density.
These are algebraic one-forms, independent of logarithm choices.
\end{proof}

Selecting $C\partial_aA/2$ afresh after a frame change need not
give the transported connection in \eqref{cbc:cf:frame-change}.
The matching theorem uses the actual connection chosen.
Where the scalar normalization has no global square root, its
rank-one module and connection retain the normalization data.

\section{Higher comparison maps}
\label{cbc:cf:cumulants-section}

The inclusion and projection are chain maps, but their products
can differ from the products of their values. Joint cumulants
organize these differences. The powers of $\hbar$ carried by
connected contractions make the resulting higher maps integral.

For a unital degree-zero linear map $F$ between graded commutative
algebras, define
\begin{equation}\label{cbc:cf:cumulants}
 \kappa_m(F)(a_1,\ldots,a_m)
  =\sum_{\pi\in\operatorname{Part}([m])}
       \epsilon(\pi)(-1)^{|\pi|-1}(|\pi|-1)!
       \prod_{B\in\pi}F\left(\prod_{i\in B}a_i\right).
\end{equation}
Within a block labels increase. Blocks are ordered by their
least label. The Koszul sign $\epsilon(\pi)$ collects the
inputs in that order. In particular
$\kappa_1(F)=F$ and
$\kappa_2(F)(a,b)=F(ab)-F(a)F(b)$.

For a graded $k[\hbar]$-module $V$, let
$S^c(V)=\bigoplus_{m\geq1}\operatorname{Sym}^m(V)$ be
the conilpotent symmetric coalgebra, with its unshuffle
coproduct. An $L_\infty[1]$ structure on $V$ is a
degree-one square-zero coderivation on $S^c(V)$.
A morphism is a coalgebra map intertwining these coderivations.
Its Taylor components are its compositions with the projection
to the cogenerators. It is a quasi-isomorphism when its unary
component is one.

Suppose $V$ is also a graded commutative algebra.
Let $Q$ be a degree-one square-zero operator of order at most
two, with $Q1=0$. Its unary part and symmetric product deviation
$B_Q$ define an $L_\infty[1]$ structure.
The second-order condition eliminates all higher operations.
In the present complexes $B_Q$ is divisible by $\hbar$.
The normalized structure has
operations $(Q,\hbar^{-1}B_Q)$.

\begin{theorem}\label{cbc:cf:higher}
For $m\geq1$,
\[
 \kappa_m(I_M)\in\hbar^{m-1}\mathcal B,\qquad
 \kappa_m(P)\in\hbar^{m-1}\mathcal B_{\mathrm{ret}}.
\]
The exact quotients
\[
 \Psi_m=\hbar^{1-m}\kappa_m(I_M),\qquad
 \Phi_m=\hbar^{1-m}\kappa_m(P)
\]
define $L_\infty[1]$ quasi-isomorphisms over $k[\hbar]$
between the normalized structures on $\mathcal B_{\mathrm{ret}}$
and $\mathcal B$. They satisfy $\Phi\circ\Psi=1$.
Their components of arity greater than one vanish on a unit
input.
\end{theorem}
\begin{proof}
In the matching formula, mark the $m$ input groups.
The partition sum \eqref{cbc:cf:cumulants} cancels every graph
disconnected on those groups. For a graph with $c$ connected
components, its coefficient is
\[
 \sum_{\sigma\in\operatorname{Part}([c])}
       (-1)^{|\sigma|-1}(|\sigma|-1)!=\delta_{c1}.
\]
This is the coefficient identity for $\log(e^t)=t$.
A connected graph on $m$ groups
has at least $m-1$ edges, each of which contributes $\hbar$
through $R_{ab}$. This proves divisibility for $I_M$.

The same argument for Wick contractions, with edge covariance
$\hbar C$, proves it for $P$. Internal edges within one
marked group are allowed. On polynomial inputs there are only
finitely many contractions.

Let $\tau_V:S^c(V)\to S^c(V)$ be the triangular coalgebra
automorphism whose Taylor components are multiplication.
Its unary component is the identity, so conilpotence gives
an inverse. The Taylor components of
\[
 \tau_{\mathrm{out}}^{-1}S^c(F)\tau_{\mathrm{in}}
\]
are the cumulants \eqref{cbc:cf:cumulants}: inversion of the
partition sum gives the coefficient
$(-1)^{|\pi|-1}(|\pi|-1)!$.
Conjugating the unary coderivation $Q$ by $\tau$ gives its
unary term and product deviation. Since $Q$ kills the unit
and has order two, all higher terms vanish.
The chain equations for $I_M$ and $P$ therefore give the
raw higher identities.

In arity $m$ each term in these identities has the known factor
$\hbar^{m-1}$. Division gives the normalized binary deviation
and the displayed higher maps. Equivalently the identities hold
after adjoining $\hbar^{-1}$ and hence in the original
polynomial modules, which have no $\hbar$-torsion.
The unary maps are quasi-isomorphisms by
Proposition~\ref{cbc:cf:matching-homotopy}.

Raw cumulants compose by cancellation of the middle $\tau$'s.
For an outer component with $q$ inputs and inner block sizes
$m_1,\ldots,m_q$, the normalization powers satisfy
\[
 (1-q)+\sum_{i=1}^q(1-m_i)=1-\sum_i m_i.
\]
Thus normalized cumulants also compose. The equality $PI_M=1$
gives $\Phi\Psi=1$. Inserting a unit in the partition sum
cancels all higher components.
\end{proof}

\begin{example}\label{cbc:cf:composite}
Use Example~\ref{cbc:cf:rotations} and abbreviate
$a_i=\eta_{y_i}$. For $\ell=2$, the two-edge matching
connects the three groups $(a_1),(a_2a_3),(a_4)$, giving
\[
 \kappa_3(I_M)(a_1,a_2a_3,a_4)
        =\hbar^2\eta_{x_1}\eta_{x_2}\eta_{x_3}\eta_{x_4}.
\]
For $\ell=3$, the connected three-edge matching gives
\[
 \Psi_4(a_1,a_2a_3,a_4a_5,a_6)
        =\eta_{x_1}\cdots\eta_{x_6}\ne0.
\]

These normalized components survive at $\hbar=0$.
For a flat metric connection the inclusion is multiplicative
and all its higher cumulants vanish. Its Wick projection may
still have higher cumulants. Even for a single Gaussian variable
of variance $\hbar/a$,
$\kappa_3(P)(x,x,x^2)=2(\hbar/a)^2$.
The third input is composite.
\end{example}

\section{Formal potential interactions}
\label{cbc:cf:interactions}

Let $V\in R[x_1,\ldots,x_d]$ and introduce a formal parameter
$g$. Work in $\mathcal B[[g]]$: each coefficient is polynomial
in $x,\hbar$ and exterior in the antifields. Define
\[
 \delta_V=g\sum_i(\partial_{x_i}V)\partial_{\eta_{x_i}}
          +g\sum_a(\partial_aV)\partial_{\eta_{y_a}}.
\]
Then $Q_{S+gV}=Q_S+\delta_V$ squares to zero.
Every perturbation raises $g$-order, so the geometric series
below are defined coefficientwise.

\begin{theorem}[Integral interacting contraction]\label{cbc:cf:interaction}
For any chosen metric connection, the maps
\begin{align}
 I_{M,g}&=(1+H_M\delta_V)^{-1}I_M,&
 P_g&=P(1+\delta_VH_M)^{-1},\notag\\
 H_{M,g}&=H_M(1+\delta_VH_M)^{-1},&
 Q_g&=Q_A+P\delta_V(1+H_M\delta_V)^{-1}I_M
 \label{cbc:cf:interacting-maps}
\end{align}
give a continuous strong deformation retract over $k[\hbar][[g]]$.
The projection is independent of $M$. After embedding in the
ring with $\hbar^{-1}$ adjoined, it is
\begin{align}
 Z_x&=P(e^{-gV/\hbar}),\notag\\
 P_g(F)&=Z_x^{-1}P(e^{-gV/\hbar}F|_{\eta_x=0}).
 \label{cbc:cf:normalized}
\end{align}
The effective differential and force are
\begin{align}
 Q_g&=\sum_a(E_a(g)-\hbar\partial_a)\partial_{\eta_{y_a}},
 \notag\\
 \sum_aE_a(g)\,dy_a
   &=dU+\frac{\hbar}{2}d\log\det A-\hbar\,d\log Z_x.
 \label{cbc:cf:interacting-force}
\end{align}
The formal interaction potential $-\hbar\log Z_x$ belongs
to $R[\hbar][[g]]$.
\end{theorem}
\begin{proof}
Apply the geometric-series contraction identity to
$(P,I_M,H_M)$. Every coefficient in $g$ uses finitely many
terms, all polynomial in $\hbar$. It gives
\[
 (Q_S+\delta_V)H_{M,g}+H_{M,g}(Q_S+\delta_V)
       =1-I_{M,g}P_g
\]
and the chain-map and side identities.

To identify $P_g$, count hidden antifields.
Each term of $H_A$ adds at least one. Its terms beyond $h_x$
add at least two, because each further $\delta_yh_x$ adds
another hidden antifield. The correction $-H_AI_MP$ in $H_M$
also adds at least two whenever it is nonzero. Indeed the
purely retained term of $I_M$ is killed by the rightmost hidden
homotopy, and every other term already has a hidden antifield.

Each $\delta_V$ removes at most one. Thus only successive
leading hidden homotopies and hidden potential derivatives can
survive the final $P$.
The projection kills inputs with hidden antifields and acts
coefficientwise on retained ones, independently of $M$.

Its functional on hidden polynomials is determined by
normalization and the hidden Ward identities. Write
$\delta_{V,x}=g\sum_iV_{x_i}\partial_{\eta_{x_i}}$.
Since $f-iPf=Q_xh_xf$, any normalized perturbed functional
$\mathcal L$ satisfies
\[
 \mathcal L(f)=Pf-\mathcal L(\delta_{V,x}h_xf).
\]
This recursion determines each successive coefficient of $g$.
The expression \eqref{cbc:cf:normalized} satisfies it by
Gaussian integration by parts, coefficientwise after adjoining
$\hbar^{-1}$. It therefore equals the integral perturbation
formula.

For a function $f$, differentiating normalized Wick evaluation
gives
\[
 \partial_aP_g(f)=P_g(\partial_af)
 -\hbar^{-1}\left(
 P_g((\partial_aS+g\partial_aV)f)
 -P_g(\partial_aS+g\partial_aV)P_g(f)\right).
\]
Together with the hidden Ward identity this proves the chain
equation on all antifield monomials. Since $P_gi=1$, the
retained differential has the displayed potential form with
$E_a(g)=P_g(\partial_aS+g\partial_aV)$.
Differentiating $Z_x$, including the varying Gaussian
normalization, gives \eqref{cbc:cf:interacting-force}.

Finally, a Wick graph with $v$ interaction vertices and $e$
edges carries $\hbar^{e-v}$. The logarithm selects connected
graphs, for which $e\geq v-1$. The outside factor $\hbar$
then gives a nonnegative power. A hidden-independent interaction
is its single classical vertex. At fixed $g$-order there are
finitely many contractions. This proves integrality of
$-\hbar\log Z_x$.
\end{proof}

\begin{corollary}\label{cbc:cf:interacting-cumulants}
For polynomial observables, including antifields,
\[
 \kappa_m(P_g)\in
       \hbar^{m-1}\mathcal B_{\mathrm{ret}}[[g]].
\]
The normalized cumulants give an $L_\infty[1]$
quasi-isomorphism for the normalized BV deviations and obey
the composition law of Theorem~\ref{cbc:cf:higher}.
\end{corollary}
\begin{proof}
Normalization cancels unmarked vacuum components. The cumulant
cancels decompositions of the marked inputs. A surviving
connected graph with $m$ marked observables and $v$ interaction
vertices has $e\geq m+v-1$, so its power $\hbar^{e-v}$
is at least $m-1$. Retained antifields supply their exterior
coefficients and Koszul signs. Moments with a hidden antifield
vanish.

The cumulant coalgebra argument then gives the higher
identities and composition before inverting $\hbar$.
The unary map is a quasi-isomorphism by
Theorem~\ref{cbc:cf:interaction}.
\end{proof}

\begin{example}[A cubic hidden interaction]\label{cbc:cf:cubic}
For one hidden coordinate, let its quadratic coefficient be
$a\in k^\times$ and let $V=yx^3/6$. Expanding
\eqref{cbc:cf:normalized} with
$P(x^{2m})=(2m-1)!!(\hbar/a)^m$ gives
\begin{align*}
 -\hbar\log Z_x
     &=-\frac{5g^2\hbar^2y^2}{24a^3}
       -\frac{5g^4\hbar^3y^4}{16a^6}+O(g^6),\\
 P_g(x)&=-\frac{g\hbar y}{2a^2}
       -\frac{5g^3\hbar^2y^3}{8a^5}+O(g^5),\\
 \hbar^{-1}\kappa_2(P_g)(x,x)
       &=\frac1a+\frac{g^2\hbar y^2}{a^4}
                  +\frac{25g^4\hbar^2y^4}{8a^7}+O(g^6),\\
 \hbar^{-2}\kappa_3(P_g)(x,x,x^2)
       &=\frac2{a^2}+\frac{8g^2\hbar y^2}{a^5}
                  +\frac{175g^4\hbar^2y^4}{4a^8}+O(g^6).
\end{align*}
The third normalized comparison is already nonzero at $g=0$
because one observable is $x^2$.
\end{example}

\section{Composition of normalized eliminations}
\label{cbc:cf:sewing}

Successive eliminations must retain the intermediate dependence
on the remaining coordinates. Let $P_x^0,P_y^0$ be independent
reference Wick evaluations, and let the interaction be $V(x,y)$.
Coefficientwise after adjoining $\hbar^{-1}$, put
\[
 Z_x=P_x^0(e^{-gV/\hbar}),\qquad
 P_x(F)=Z_x^{-1}P_x^0(e^{-gV/\hbar}F),\qquad
 P_y^{\mathrm{eff}}(G)
       =\frac{P_y^0(Z_xG)}{P_y^0(Z_x)}.
\]
The factor $Z_x$ cancels in the composite, so
\[
 P_y^{\mathrm{eff}}P_x(F)=
 \frac{P_y^0P_x^0(e^{-gV/\hbar}F)}
      {P_y^0P_x^0(e^{-gV/\hbar})}.
\]

This is joint normalized evaluation. The equality descends to
the integral formal maps above wherever the successive
nondegenerate Gaussian contractions are defined. Their
normalized cumulants compose by
Theorem~\ref{cbc:cf:higher}.

For a nondegenerate Gaussian block matrix, the same identity
is expressed by completing the square. If
\[
 K=\begin{pmatrix}A&B\\B^{\mathsf t}&D\end{pmatrix},
 \qquad S_D=D-B^{\mathsf t}A^{-1}B,
\]
with $A$ and $S_D$ invertible in the specified coefficient
ring, elimination of $x$ leaves the quadratic form $S_D$.
The triangular change
$x\mapsto x+A^{-1}By$ has determinant one, and
\[
 \det K=(\det A)(\det S_D).
\]
It therefore composes both the Wick contractions and their
normalization lines.

For an interacting example, take reference variances
$P_x^0(x^2)=P_y^0(y^2)=\hbar$ and $V=yx^3/6$. Then
\[
 Z_x=1+\frac5{24}g^2\hbar y^2+O(g^4),\qquad
 P_y^{\mathrm{eff}}(y^2)
       =\hbar+\frac5{12}g^2\hbar^3+O(g^4).
\]
The correction uses
$P_y^0(y^4)-P_y^0(y^2)^2=2\hbar^2$.
Replacing the intermediate factor by a scalar would remove it.

These constructions determine maps on their polynomial or formal
coefficient modules. The line gauge additionally retains all
position pairs and their logarithmic residue. To compare them
with the ordered collision coefficients of
Chapters~\ref{cbc:chap:heisenberg} and~\ref{cbc:chap:transfer},
one must construct the maps on those coefficient systems.

\begin{problem}[The chiral comparison]\label{cbc:cf:chiral-comparison}
Construct a comparison from the transverse collision coefficients
and their varying endpoint modules to the all-mode line
connection. Require it to preserve the polynomial spectral maps.
Specify the continuous coefficient dual, the maps through
collision strata and coherent gauges on nested collision
domains. For a BV realization, also construct the supported
operation maps, state domains and renormalized contractions
whose finite reductions are the Gaussian maps above.
Require compatibility with both vacuum insertions and with
composition of normalized eliminations.
\end{problem}
